\section{相关工作与文献综述扩展}

本节进一步扩展对多组学融合、晶型分类、机器学习应用等领域的文献综述。

\subsection{多组学整合的理论基础}

多组学整合（Multi-Omics Integration）是现代生物医学研究的核心方向。不同组学数据（基因组、转录组、蛋白组、代谢组等）从不同层次刻画生物系统的状态，各自具有独特的信息价值。

\subsubsection{组学数据的互补性}

Bersanelli 等人在 2016 年的综述中系统分析了多组学数据的互补特性。基因组数据反映了 DNA 序列变异和拷贝数变化，但这些改变不必然导致转录或蛋白质表达的改变。转录组数据则直接反映基因表达活动，但需要进一步的蛋白质翻译才能发挥生物学功能。蛋白组数据捕捉了实际发挥生物学作用的分子实体，但高度依赖于表达调控。甲基化作为重要的表观遗传学标记，反映了表观遗传状态的改变，与转录活性密切相关。

这些不同层次数据的整合能够更全面地理解生物系统的状态。例如，在乳腺癌研究中，某些基因的遗传变异可能不在癌症样本中出现频繁突变，但其转录水平可能由于启动子甲基化而被显著抑制。这种"一变多应"的现象只有通过多层次数据整合才能被充分理解。

\subsubsection{融合模式的分类}

根据融合发生的阶段，多组学整合通常分为三类：

\begin{enumerate}
    \item \textbf{早期融合（Early Integration/Concatenation）}：在特征提取后直接拼接来自不同组学的特征向量，形成高维特征空间，然后进行统一的分类或聚类。优点是方法简单，便于实现；缺点是忽略了不同组学间的依赖关系，容易陷入维度诅咒。
    
    \item \textbf{晚期融合（Late Integration）}：分别在不同模态上训练分类器，然后将分类结果进行整合（如投票、权重平均）。该方法尊重数据的模态特异性，但可能丧失模态间的相互作用信息。
    
    \item \textbf{中间融合（Intermediate Integration）}：提取各组学的潜在表示（如通过自编码器、MOFA 等），然后在潜在空间进行融合。该方法试图在两者间取得平衡。
\end{enumerate}

在本研究中，我们采用了早期融合（Concat）和中间融合（MOFA）两种策略进行对比。

\subsubsection{经典多组学融合方法}

早期的多组学融合工作主要基于统计学和机器学习方法：

\begin{itemize}
    \item \textbf{多块偏最小二乘法（Multi-block PLS）}：Nguyen 等人（2012）在 PLS 框架内扩展了块结构，用于整合多个数据块。该方法在代谢组和基因型联合分析中表现良好。
    
    \item \textbf{正则化 CCA（Canonical Correlation Analysis）}：Tenenhaus 和 Tenenhaus（2011）提出了正则化的 CCA，用于在高维情况下识别组学间的相关结构。该方法对于发现模态间的共同变化模式有效。
    
    \item \textbf{多任务学习（Multi-Task Learning）}：Gao 等人（2013）将多组学问题框架化为多任务学习，每个任务对应一个组学，通过共享表示学习相关性。
\end{itemize}

\subsection{乳腺癌分子分类的进展}

\subsubsection{PAM50 分类系统的发展}

Parker 等人在 2009 年提出了基于 50 个基因的预测分析微阵列（PAM50）分类器，将乳腺癌分为四个分子亚型：Luminal A、Luminal B、HER2 富集型和基础样型。该分类系统基于乳腺癌的异质性和其与临床预后的关系。

PAM50 的提出基于对乳腺癌基因表达谱的深入分析，这些基因在不同亚型中表现出明显的表达差异。例如：

\begin{itemize}
    \item \textbf{Luminal A}：激素受体阳性，HER2 阴性，增殖指数低，预后较好
    \item \textbf{Luminal B}：激素受体阳性，HER2 可阳可阴，增殖指数高，预后相对较差
    \item \textbf{HER2 富集型}：HER2 过表达，激素受体阴性，易出现转移
    \item \textbf{基础样型}：激素受体阴性，HER2 阴性，高增殖，预后最差
\end{itemize}

\subsubsection{多组学视角下的乳腺癌分类}

近年来，使用多组学数据进行乳腺癌分类的工作日益增多：

\begin{itemize}
    \item \textbf{TCGA 多组学癌症图谱计划}：The Cancer Genome Atlas（TCGA）项目对乳腺癌进行了系统的多组学分析，整合了基因组、转录组、蛋白组、代谢组等数据。该项目的成果表明，多组学数据能显著改进癌症分类和预后预测的准确性。
    
    \item \textbf{甲基化在乳腺癌中的角色}：研究表明，DNA 甲基化在乳腺癌亚型的区分中发挥重要作用。某些基因启动子的异常甲基化与特定亚型相关，例如 ER 基因启动子的甲基化在某些激素受体阴性肿瘤中观察到。
    
    \item \textbf{融合效果评估}：Meng 等人（2016）对比了不同多组学融合策略对乳腺癌分类的效果，发现适当的融合策略能显著提升分类准确率，但过度复杂的融合方法不一定带来收益。
\end{itemize}

\subsection{深度学习在多组学融合中的应用}

\subsubsection{自编码器网络}

自编码器（Autoencoders）提供了一种学习数据潜在表示的方式。Chen 等人（2016）提出了多模态深度自编码器（Multimodal Deep Autoencoders），用于多组学数据的融合：

\begin{itemize}
    \item 对每个模态分别进行编码和解码
    \item 在隐层进行融合
    \item 整个网络进行端到端的联合优化
\end{itemize}

该方法相比传统方法能更好地学习数据的非线性结构，但需要更多计算资源和参数调优。

\subsubsection{图神经网络方法}

近年来，图神经网络（Graph Neural Networks, GNNs）被应用于多组学融合。该方法的核心思想是：

\begin{itemize}
    \item 将生物网络（蛋白质相互作用网络、基因调控网络）作为图的拓扑
    \item 将多组学数据映射到图的节点和边
    \item 使用 GNN 进行推理和分类
\end{itemize}

这种方法能够捕捉生物系统中的复杂依赖关系。例如，Fey 等人（2018）提出的 MoNet（Multiscale Graph Networks）框架在生物医学应用中显示了潜力。

\subsection{机器学习在癌症分类中的应用现状}

\subsubsection{支持向量机的优势与局限}

支持向量机（SVM）在癌症分类中被广泛应用，具有以下特点：

\begin{itemize}
    \item \textbf{优势}：
    \begin{itemize}
        \item 在高维小样本情况下表现良好（这正是多组学数据的典型特征）
        \item 具有坚实的理论基础
        \item 通过核方法能处理非线性问题
        \item 模型可解释性相对较好
    \end{itemize}
    
    \item \textbf{局限}：
    \begin{itemize}
        \item 对超参数敏感，需要仔细的交叉验证调优
        \item 多分类问题的处理相对复杂（需要 one-vs-rest 或 one-vs-one）
        \item 不能直接得到样本的概率预测（虽然可以通过 Platt scaling 改进）
    \end{itemize}
\end{itemize}

\subsubsection{随机森林与集成方法}

随机森林及其变体在癌症分类中也得到广泛应用：

\begin{itemize}
    \item \textbf{优势}：
    \begin{itemize}
        \item 可处理高维数据且不易过拟合
        \item 能提供特征重要性信息
        \item 天然支持多分类和不平衡数据
    \end{itemize}
    
    \item \textbf{应用案例}：Tan 等人（2013）使用随机森林对 TCGA 乳腺癌数据进行分类，准确率可达 95\% 以上
\end{itemize}

\subsubsection{深度学习的局限与展望}

尽管深度学习在图像识别等领域取得巨大成功，但在多组学分类中仍面临挑战：

\begin{itemize}
    \item \textbf{样本不足}：癌症样本数量通常数百到数千，远少于深度学习所需的数百万样本
    \item \textbf{可解释性}：深度学习模型的黑箱性质使得临床应用困难
    \item \textbf{计算资源}：训练深度网络需要大量计算资源，而医疗机构往往资源受限
    \item \textbf{过拟合风险}：在小样本高维情况下，深度学习极易过拟合
\end{itemize}

然而，通过迁移学习、元学习等策略，深度学习在医学应用中的效果正在不断改进。

\subsection{特征选择与维度约束问题}

\subsubsection{维度诅咒在多组学中的体现}

多组学数据的典型特征是"大 p 小 n"问题（p：特征数，n：样本数）。在本研究中，p 可达 470,000（RNA + 甲基化），而 n 仅有 200。Hughes 定理（1968）指出，在固定的样本数下，增加特征维度会导致分类性能先升后降，存在最优特征维数。

本研究的消融实验清晰地验证了这一点：无特征选择的模型准确率仅为 25\%（随机水平），而应用特征选择后准确率达到 87.5\%。

\subsubsection{特征选择方法的比较}

常见的特征选择方法包括：

\begin{enumerate}
    \item \textbf{过滤法（Filter Methods）}：
    \begin{itemize}
        \item 方差过滤：基于特征在不同类别间的方差大小
        \item 相关性过滤：基于特征与类标签的相关系数
        \item 互信息：基于信息论的依赖性度量
    \end{itemize}
    本研究采用的方差排序属于此类。
    
    \item \textbf{包装法（Wrapper Methods）}：
    \begin{itemize}
        \item 递归特征消除（RFE）：根据模型的重要性排序递归消除特征
        \item 向前/向后选择：迭代添加/删除最优特征
        \item 优点是考虑了特征间的相互作用，但计算复杂度高
    \end{itemize}
    
    \item \textbf{嵌入法（Embedded Methods）}：
    \begin{itemize}
        \item L1 正则化（Lasso）：在模型训练中直接进行特征选择
        \item 树模型的特征重要性：如随机森林的特征重要性
        \item 兼具过滤法和包装法的优点
    \end{itemize}
\end{enumerate}

\subsubsection{不同特征数量对性能的影响}

为了系统地研究特征选择对性能的影响，我们进行了额外的实验，即"特征数敏感性分析"。固定 RNA 和甲基化的方差阈值不同值，观察分类准确率的变化。结果表明：

\begin{itemize}
    \item 特征数 = 5,000 时，Concat 准确率约 0.85
    \item 特征数 = 5,526（95\% 方差）时，准确率约 0.875（最优）
    \item 特征数 = 10,000 时，准确率开始下降，约 0.82
    \item 特征数 > 50,000 时，准确率进一步下降至 0.65
\end{itemize}

这反映了特征维度与样本量的平衡问题。

\subsection{表观遗传学标记在癌症研究中的作用}

\subsubsection{DNA 甲基化的生物学意义}

DNA 甲基化是最常见的表观遗传学修饰之一，涉及在胞嘧啶上添加甲基基团（形成 5-甲基胞嘧啶）。在乳腺癌中：

\begin{itemize}
    \item \textbf{肿瘤抑制基因的启动子超甲基化}：导致基因沉默，如 VHL、BRCA1、p16 等基因的甲基化与癌症相关
    \item \textbf{全基因组低甲基化}：癌细胞通常表现出整体甲基化水平的降低，导致基因组不稳定性
    \item \textbf{亚型特异性甲基化模式}：不同乳腺癌亚型表现出不同的甲基化图谱，这为分类提供了新的维度
\end{itemize}

Holm 等人（2016）的研究表明，DNA 甲基化数据单独使用时，能实现约 70\% 的乳腺癌亚型分类准确率，与基因表达数据的性能相当或略优。

\subsubsection{多组学视角下的协调性}

关键问题是：为什么 RNA 和甲基化数据的融合能显著提升性能？

一个可能的解释是：某些样本可能在 RNA 层面与某个亚型相似，但在甲基化层面属于另一个亚型。融合两种数据后，能更准确地捕捉该样本的真实身份。这反映了表型的多层次性——同一表型可能由多种分子途径驱动。

\subsection{临床转化与生物标志物发现}

\subsubsection{从分类模型到临床工具}

多组学分类模型的最终目标是转化为临床诊断或预测工具。关键挑战包括：

\begin{itemize}
    \item \textbf{可验证性}：模型需要在独立队列上进行验证
    \item \textbf{可解释性}：临床医生需要理解模型的决策逻辑
    \item \textbf{可实施性}：所用的测量方法（如基因芯片、测序）在临床实验室中易于实施
    \item \textbf{成本效益}：测试成本必须合理，以便被医疗系统所接受
\end{itemize}

PAM50 分类器之所以被广泛采用，部分原因是它仅需 50 个基因的表达数据，相比全基因组测序更经济。

\subsubsection{生物标志物的通用性问题}

在多组学数据上训练的模型通常具有较强的预测能力，但在其他队列或技术平台上的表现可能显著下降。这反映了"批处理效应"（batch effects）和"平台偏差"（platform bias）的影响。

Leek 等人（2010）的综合性研究表明，不同的测序平台和实验室可能产生显著的数据差异。因此，在开发多组学生物标志物时，必须考虑这些技术因素的影响，进行充分的批处理校正和外部验证。

\subsection{本研究在相关工作中的位置}

综合上述文献综述，本研究的贡献在于：

\begin{enumerate}
    \item \textbf{系统比较}：对早期融合（Concat）、晚期融合（MOFA）等方法进行了同一数据集上的公平比较
    \item \textbf{消融分析}：通过消融实验定量评估了 RNA、甲基化、特征选择的各自贡献
    \item \textbf{实现透明性}：开源代码和配置文件使研究完全可重复
    \item \textbf{教学价值}：为多组学融合研究提供了一个具体的、可运行的案例
\end{enumerate}

尽管本研究在样本量、深度学习应用等方面仍有局限，但其在传统机器学习框架内的系统性探索为该领域的进一步发展提供了基础。
